\documentclass[12pt]{article}
\usepackage[utf8]{inputenc}
\usepackage[T1]{fontenc}
\usepackage{hyperref}
\title{Ukit Language Resources for K-12 Curriculum Development}
\author{Compiled by Claude Sonnet 5 (claude-sonnet-5) for Andr\'as Kornai}
\date{}
\begin{document}
\maketitle

\section{Introduction}
Ukit (ISO 639-3 umi) is a Punan language of Sarawak, Malaysia. It is described as endangered and unwritten. Speaker counts conflict (about 120 in a 1981 figure, about 200 from Ethnologue 2016 via Joshua Project). The Ukit people live mainly in the upper Rajom and Tatau river areas of the Baleh region. Do not confuse Ukit with Bukitan (bkn), which is mainly in West Kalimantan.

\section{Primary Educational Resources}
The search found none. The language is not known to be taught in schools.

\section{Linguistic and Reference Materials}
\begin{itemize}
\item Wikipedia article ``Ukit language'': \url{https://en.wikipedia.org/wiki/Ukit_language}
\item Wikipedia article ``Ukit people'': \url{https://en.wikipedia.org/wiki/Ukit_people}
\item Wikipedia article ``Bukitan language'' (a related but distinct language): \url{https://en.wikipedia.org/wiki/Bukitan_language}
\end{itemize}

\section{Religious and Community Resources}
Joshua Project profile of the Ukit in Malaysia: \url{https://joshuaproject.net/people_groups/15694/MY}

\section{Digital and Audio Resources}
None found.

\section{Curriculum Development Context}
The search found almost nothing beyond encyclopedic summaries. Fieldwork, community consultation and orthography development would have to come first. No community organizations were found.
\end{document}